\documentclass[11pt,a4paper]{article}
\usepackage[a4paper,top=5mm,bottom=5mm,left=13mm,right=13mm]{geometry}
\usepackage[T1]{fontenc}
\usepackage{amsmath}
\usepackage{amssymb}
\usepackage{enumitem}
\usepackage{multicol}
\usepackage{tikz}
\usetikzlibrary{decorations.pathreplacing}

\pagestyle{empty}
\setlength{\parindent}{0pt}
\setlist[enumerate,1]{label=\textbf{Q\arabic*.}, leftmargin=*, itemsep=6pt, topsep=4pt, parsep=0pt}
\setlist[enumerate,2]{label=(\alph*), leftmargin=*, itemsep=2pt, topsep=1pt, parsep=0pt}

\begin{document}

\begin{center}
  {\Large\bfseries Percentages: Non-calculator questions --- Worked Solutions}
\end{center}

\begin{enumerate}

%----------------------------------------------------------------
\item A badge has a grid of \(100\) squares, and \(51\) squares are coloured.
\begin{enumerate}
  \item Since there are \(100\) squares, the percentage coloured is the number of coloured squares out of \(100\):
  \[
  \frac{51}{100}=51\%.
  \]
  \item \(51\%=\dfrac{51}{100}=0.51\).
  \item \(51\%\) as a fraction of the badge is \(\dfrac{51}{100}\).
\end{enumerate}

%----------------------------------------------------------------
\item Fill in the gaps.  Multiply by \(100\) to change a decimal to a percentage; divide by \(100\) to change a percentage to a decimal.
\begin{multicols}{3}
\begin{enumerate}
  \item \(0.4=40\%\)
  \item \(0.06=6\%\)
  \item \(1.25=125\%\)
  \item \(45\%=0.45\)
  \item \(7\%=0.07\)
  \item \(130\%=1.3\)
\end{enumerate}
\end{multicols}

%----------------------------------------------------------------
\item From the bar chart, the numbers of votes are:
\[
\text{Neon}=18,\qquad \text{Retro}=12,\qquad \text{Space}=15,\qquad \text{Jungle}=5.
\]
There are \(50\) votes altogether.
\begin{enumerate}
  \item Neon:
  \[
  \frac{18}{50}\times 100\%=36\%.
  \]
  \item Jungle:
  \[
  \frac{5}{50}\times 100\%=10\%.
  \]
  \item Space or Retro received
  \[
  15+12=27
  \]
  votes out of \(50\), which is
  \[
  \frac{27}{50}\times 100\%=54\%.
  \]
  Since \(54\%>50\%\), Ravi is right.
\end{enumerate}

%----------------------------------------------------------------
\item Work these out:
\begin{multicols}{3}
\begin{enumerate}
  \item \(10\%\) of \(250=\frac{1}{10}\times 250=25\).
  \item \(10\%\) of \(60=6\), so \(5\%\) of \(60=3\).
  \item \(10\%\) of \(80=8\), so \(30\%=24\) and \(5\%=4\). Hence
  \[
  35\%\text{ of }80=24+4=28.
  \]
\end{enumerate}
\end{multicols}

%----------------------------------------------------------------
\item Work out each new price.
\begin{multicols}{2}
\begin{enumerate}
  \item \(10\%\) of \pounds45 is \pounds4.50, so \(20\%\) is \pounds9.
  The trainers are reduced by \pounds9, so the sale price is \pounds36.

  \item \(25\%\) of \pounds32 is \pounds8.
  The hoodie is increased by \pounds8, so the new price is \pounds40.
\end{enumerate}
\end{multicols}

%----------------------------------------------------------------
\item
\begin{enumerate}
  \item A reduction of \(20\%\) means the coat now costs \(80\%\) of its original price.
  \[
  80\%=\pounds48
  \]
  so
  \[
  10\%=\pounds6
  \]
  and therefore
  \[
  100\%=\pounds60.
  \]
  The original price was \pounds60.

  \item After a \(30\%\) increase, the new price is \(130\%\) of the original price.
  \[
  130\%=\pounds91
  \]
  so
  \[
  10\%=91\div 13=\pounds7
  \]
  and therefore
  \[
  100\%=7\times 10=\pounds70.
  \]
  The membership before the rise was \pounds70 per month.
\end{enumerate}

%----------------------------------------------------------------
\item For a repair cost of \pounds64:
\[
10\%\text{ of }64=6.40,
\]
so
\[
5\%=6.40\div 2=3.20
\]
and
\[
2.5\%=3.20\div 2=1.60.
\]
Therefore
\[
17.5\%=10\%+5\%+2.5\%=6.40+3.20+1.60=11.20.
\]
The VAT is \pounds11.20.

%----------------------------------------------------------------
\item Coffee and tea together make
\[
35\%+20\%=55\%.
\]
So the cold drinks are
\[
100\%-55\%=45\%.
\]
Now
\[
10\%\text{ of }240=24,
\]
so
\[
40\%\text{ of }240=96
\]
and
\[
5\%\text{ of }240=12.
\]
Hence
\[
45\%\text{ of }240=96+12=108.
\]
The caf\'e sold \(108\) cold drinks.

%----------------------------------------------------------------
\item The customer is \textbf{not} right.

Take an example price of \pounds100.
\begin{itemize}
  \item After a \(20\%\) increase:
  \[
  100\times 1.20=120.
  \]
  \item Then \(20\%\) is taken off the new, larger price:
  \[
  20\%\text{ of }120=24.
  \]
  \item Final sale price:
  \[
  120-24=96.
  \]
\end{itemize}
So the final price is \pounds96, not \pounds100. The \(20\%\) decrease takes off \pounds24, while the original \(20\%\) increase had added only \pounds20.

%----------------------------------------------------------------
\item The normal price of the headphones is \pounds80.
\begin{enumerate}
  \item Shop A: \(30\%\) off \pounds80 means the customer pays
  \[
  70\%\text{ of }80=0.70\times 80=56,
  \]
  so the price is \pounds56.

  Shop B: a quarter off means the customer pays
  \[
  75\%\text{ of }80=0.75\times 80=60,
  \]
  and then a further \pounds8 off gives
  \[
  60-8=52.
  \]
  So Shop B is cheaper, by
  \[
  \pounds56-\pounds52=\pounds4.
  \]

  \item One possible third offer is: take off \(20\%\), then take off a further \pounds15.

  For \pounds80:
  \[
  20\%\text{ of }80=16,
  \]
  so the price after \(20\%\) off is
  \[
  80-16=64.
  \]
  Taking off another \pounds15 gives
  \[
  64-15=49.
  \]
  This is cheaper than Shop A's \pounds56 and Shop B's \pounds52. Also, half of the normal price is \pounds40, and \pounds49 is more than \pounds40.

  \item Let the normal price be \(x\) pounds.
  \[
  \text{Shop A price}=0.70x.
  \]
  \[
  \text{Shop B price}=0.75x-8.
  \]
  They charge the same when
  \[
  0.70x=0.75x-8.
  \]
  Subtract \(0.70x\) from both sides:
  \[
  0=0.05x-8
  \]
  so
  \[
  0.05x=8
  \]
  and
  \[
  x=\frac{8}{0.05}=160.
  \]
  The headphones would need to cost \pounds160 for both shops to charge the same price. Both would then charge \pounds112.
\end{enumerate}

\end{enumerate}

\end{document}